\chapter{Counterexamples}
\label{chap:counterexamples}

Counterexamples that separate notions which agree in finite dimensions. This chapter collects
those about immersions and Finsler structures: a weak immersion that is not an immersion
($\ell^1 \to \ell^2$), an injective immersion that is not initial (the figure eight), an
immersion between Banach spaces that is not split ($c_0 \subseteq \ell^\infty$, by Phillips'
lemma), and a Finsler length structure that is not a Finsler metric. Three more are stated next
to the results they limit: McAlpin's complete Riemannian manifold without shortest curves
(Example~\ref{ex:ms-mcalpin}), a bounded Finsler structure on an incomplete model whose distance
does not give the manifold topology (Theorem~\ref{thm:ms-incomplete-counterexample}), and a weak
Riemannian metric without Levi-Civita connection (Theorem~\ref{thm:mt-no-levi-civita}).

\section{Weak immersions, injective immersions and split immersions}

The notions of immersion form a chain: split immersions are immersions,
immersions are weak immersions, and injective immersions need not be initial. In finite
dimensions the first two notions coincide. The three examples of this section show that every
implication is strict, and that the first two separations occur already between Banach spaces.

\subsection{A weak immersion that is not an immersion}

Let $\iota \colon \ell^1 \to \ell^2$ be the inclusion of real sequences. Both spaces are Banach
spaces, hence manifolds with a single chart.

\begin{proposition}[$\ell^1 \to \ell^2$]
  \label{prop:ex-l1-l2}
  \lean{LocallyConvex.Manifolds.lpInclusion, LocallyConvex.Manifolds.lpInclusion_injective, LocallyConvex.Manifolds.not_boundedBelow_lpInclusion, LocallyConvex.Manifolds.isWeakImmersion_lpInclusion, LocallyConvex.Manifolds.not_isImmersion_lpInclusion}
  \leanok
  \uses{def:weak-immersion, def:immersion}
  The inclusion $\iota$ is continuous, linear and injective, hence a weak immersion. It is not
  bounded below, hence it is not an immersion.
\end{proposition}

\begin{proof}
  \leanok
  Continuity follows from the closed graph theorem: the graph of $\iota$ is closed because the
  coordinate evaluations are continuous on both spaces. The differential of $\iota$ at every
  point is $\iota$, which is injective. For $x_n = e_0 + \dots + e_{n-1}$ we have
  $\|x_n\|_1 = n$ and $\|x_n\|_2 = \sqrt n$, so no inequality $m\|x\|_1 \le \|\iota x\|_2$ with
  $m > 0$ can hold, and $\iota$ is not a topological isomorphism onto its range.
\end{proof}

\begin{remark}
  The range of $\iota$ is dense in $\ell^2$ and not closed. In finite dimensions an injective
  linear map is automatically a topological isomorphism onto a closed subspace; in infinite
  dimensions the condition that the differential be an embedding with closed range is a genuine
  additional requirement, and this is the reason to distinguish weak immersions from
  immersions.
\end{remark}

\subsection{An injective immersion that is not initial: the figure eight}

Let $\theta(t) = 2\arctan t \in (-\pi, \pi)$ and
\[
  j \colon \R \to \R^2, \qquad j(t) = (\sin\theta(t), \sin 2\theta(t)).
\]
As $t$ runs over $\R$, $j(t)$ traces the figure eight (a lemniscate), passing through the origin
at $t = 0$ and tending to the origin again as $t \to \pm\infty$.

\begin{proposition}[The figure eight is immersed]
  \label{prop:ex-lemniscate-immersed}
  \lean{LocallyConvex.Manifolds.Lemniscate.j, LocallyConvex.Manifolds.Lemniscate.j_injective, LocallyConvex.Manifolds.Lemniscate.vel_ne_zero, LocallyConvex.Manifolds.Lemniscate.isImmersedSubmanifold_lemniscate}
  \leanok
  \uses{def:immersion}
  The map $j$ is smooth, injective and has nonvanishing velocity; hence it is an injective
  immersion.
\end{proposition}

\begin{proof}
  \leanok
  On $(-\pi, \pi)$ the pair $(\sin a, \sin 2a)$ determines $a$, which gives injectivity. If
  $\dot j(t) = 0$, then $\cos\theta = 0$ and $\cos 2\theta = 0$ at the same time, which is
  impossible. A nonzero vector spans a closed line, so $v \mapsto v\,\dot j(t)$ is a
  topological isomorphism onto a closed subspace.
\end{proof}

\begin{theorem}[The figure eight is not initial]
  \label{thm:ex-lemniscate-not-initial}
  \lean{LocallyConvex.Manifolds.Lemniscate.f, LocallyConvex.Manifolds.Lemniscate.g, LocallyConvex.Manifolds.Lemniscate.j_g, LocallyConvex.Manifolds.Lemniscate.one_lt_g, LocallyConvex.Manifolds.Lemniscate.g_zero, LocallyConvex.Manifolds.Lemniscate.not_isInitialSubmanifold_lemniscate}
  \leanok
  \uses{prop:ex-lemniscate-immersed, def:initial-submanifold}
  Let $\sigma(s) = \tfrac12\arctan s$ and $f(s) = (\sin\sigma, -\sin 2\sigma)$. Then $f$ is a
  smooth curve in $\R^2$ with values in $j(\R)$, but its unique lift $g$ with $j \circ g = f$ is
  not continuous at $0$. Hence $j$ is not an initial submanifold.
\end{theorem}

\begin{proof}
  \leanok
  The curve $f$ crosses the origin along the branch that $j$ reaches only as $t \to \pm\infty$.
  Explicitly, $g = \tan(\lambda/2)$ with $\lambda(s) = \pi - \sigma(s)$ for $s > 0$,
  $\lambda(s) = -\pi - \sigma(s)$ for $s < 0$ and $\lambda(0) = 0$; trigonometric identities give
  $j \circ g = f$. For $s > 0$ we have $g(s) > 1$, while $g(0) = 0$. If $j$ were initial, $g$
  would be smooth, in particular continuous.
\end{proof}

\begin{remark}
  This is the classical finite dimensional example; it shows that the condition defining
  initial submanifolds is not a consequence of being an injective immersion, and has nothing to
  do with infinite dimensions. The topology of $\R$ transported by $j$ is strictly finer than the
  topology induced by $\R^2$ on the figure eight.
\end{remark}

\subsection{An immersion that is not split: \texorpdfstring{$c_0 \subseteq \ell^\infty$}{c0 in l-infinity}}

We realize $\ell^\infty$ as the Banach space of bounded real sequences and $c_0$ as the closed
subspace of sequences tending to $0$. The classical obstruction is the following result of
Phillips \cite[7.5, pp.~539--540]{Phi40}, stated there for the space $c$ of convergent
sequences, which contains $c_0$ with codimension one; we formalize the short proof of Whitley \cite{Whi66}.

\begin{theorem}[Phillips' lemma]
  \label{thm:ex-phillips}
  \lean{LocallyConvex.Manifolds.Phillips.Linf, LocallyConvex.Manifolds.Phillips.c0, LocallyConvex.Manifolds.Phillips.isClosed_c0, LocallyConvex.Manifolds.Phillips.not_closedComplemented_c0}
  \leanok
  There is no continuous linear projection of $\ell^\infty$ onto $c_0$.
\end{theorem}

The proof rests on an uncountable family of almost disjoint subsets of $\N$, and on the fact
that a functional vanishing on $c_0$ can only see countably many of them.

\begin{lemma}[An almost disjoint family]
  \label{lem:ex-almost-disjoint}
  \lean{LocallyConvex.Manifolds.Phillips.N, LocallyConvex.Manifolds.Phillips.infinite_N, LocallyConvex.Manifolds.Phillips.finite_N_inter, LocallyConvex.Manifolds.Phillips.ind, LocallyConvex.Manifolds.Phillips.ind_not_mem_c0}
  \leanok
  For each $A \subseteq \N$ there is an infinite set $N_A \subseteq \N$ such that
  $N_A \cap N_B$ is finite whenever $A \ne B$. The indicator $\chi_A \in \ell^\infty$ of $N_A$
  does not belong to $c_0$.
\end{lemma}

\begin{proof}
  \leanok
  Index the finite binary words by $\N$, and let $N_A$ be the set of indices of the initial
  segments of the indicator function of $A$. If $A$ and $B$ differ at $i_0$, their initial
  segments of length larger than $i_0$ are different, so $N_A \cap N_B$ is finite. Since $N_A$
  is infinite, $\chi_A$ takes the value $1$ infinitely often.
\end{proof}

\begin{lemma}[Whitley]
  \label{lem:ex-whitley}
  \lean{LocallyConvex.Manifolds.Phillips.abs_le_of_eventually, LocallyConvex.Manifolds.Phillips.countable_ne_zero}
  \leanok
  \uses{lem:ex-almost-disjoint}
  Let $f$ be a continuous linear functional on $\ell^\infty$ that vanishes on $c_0$. If
  $|x_k| \le 1$ for all but finitely many $k$, then $|f(x)| \le \|f\|$. Consequently,
  $f(\chi_A) \ne 0$ only for countably many $A \subseteq \N$.
\end{lemma}

\begin{proof}
  \leanok
  For the first claim, modify $x$ on a finite set to obtain $y$ with $\|y\| \le 1$; then
  $x - y \in c_0$ and $f(x) = f(y)$. For the second, suppose $|f(\chi_A)| > 1/m$ for every $A$ in
  a finite set $\mathcal F$, and let $x = \sum_{A \in \mathcal F} \pm\chi_A$, with the signs of
  $f(\chi_A)$. Off the finite set where two of the $N_A$ meet, $|x_k| \le 1$; hence
  $\#\mathcal F/m < f(x) \le \|f\|$. So each set $\{A : |f(\chi_A)| > 1/m\}$ is finite.
\end{proof}

\begin{proof}[Proof of Theorem~\ref{thm:ex-phillips}]
  \proves{thm:ex-phillips}
  \leanok
  \uses{lem:ex-whitley}
  Suppose $P$ is a continuous projection onto $c_0$, and apply Lemma~\ref{lem:ex-whitley} to the
  functionals $f_n = e_n \circ (1 - P)$, where $e_n$ is the $n$-th coordinate. The subsets of
  $\N$ are uncountable, so some $A$ avoids the countable union of the exceptional sets. Then
  $(1 - P)\chi_A = 0$, that is, $\chi_A = P\chi_A \in c_0$, a contradiction.
\end{proof}

\begin{corollary}[Embedded but not split]
  \label{cor:ex-c0}
  \lean{LocallyConvex.Manifolds.Phillips.isEmbeddedSubmanifold_c0, LocallyConvex.Manifolds.Phillips.not_isSplitSubmanifold_c0, LocallyConvex.Manifolds.Phillips.isImmersion_c0, LocallyConvex.Manifolds.Phillips.not_isSplitImmersion_c0}
  \leanok
  \uses{thm:ex-phillips, prop:ex-subspace, prop:ex-subspace-immersion,
    def:embedded-submanifold, def:split-submanifold, def:immersion, def:split-immersion}
  The subspace $c_0$ is an embedded submanifold of $\ell^\infty$ that is not split, and the
  inclusion $c_0 \to \ell^\infty$ is an immersion that is not a split immersion.
\end{corollary}

\begin{proof}
  \leanok
  Combine Propositions~\ref{prop:ex-subspace} and~\ref{prop:ex-subspace-immersion} with
  Theorem~\ref{thm:ex-phillips}.
\end{proof}

\begin{remark}
  In Hilbert spaces every closed subspace is complemented, so the distinction between embedded
  and split submanifolds, and between immersions and split immersions, is invisible there. For
  Banach spaces it is not, and Corollary~\ref{cor:ex-c0} gives the standard witness.
\end{remark}

\section{A Finsler length structure that is not a Finsler metric}
\label{sec:mt-cusp}

The difference between the two notions is not an infinite-dimensional phenomenon. On $\R^2$ with
the sup norm, let
\[
  \beta(x, y) = \frac{\max(0, x^2 - |y|)}{x^2 + |y|} \quad (\beta(0,0) = 0), \qquad
  \mu_p(v) = \|v\| + \beta(p)\,|v_2| .
\]
The weight $\beta$ takes values in $[0,1]$, is supported in the cusp $\{|y| < x^2\}$, equals $1$
on the $x$-axis minus the origin, and is continuous away from the origin.

\begin{proposition}
  \label{prop:mt-cusp-fibers}
  \lean{LocallyConvex.Manifolds.CuspFinsler.field, LocallyConvex.Manifolds.CuspFinsler.isNondegenerate, LocallyConvex.Manifolds.CuspFinsler.isReversible, LocallyConvex.Manifolds.CuspFinsler.continuousAt_β}
  \leanok
  \uses{def:mt-field}
  Each $\mu_p$ is a reversible Finsler norm.
\end{proposition}

\begin{proof}
  \leanok
  It is the sum of the norm and a nonnegative multiple of the seminorm $|v_2|$.
\end{proof}

\begin{theorem}
  \label{thm:mt-cusp}
  \lean{LocallyConvex.Manifolds.CuspFinsler.not_continuous, LocallyConvex.Manifolds.CuspFinsler.lengthStructure, LocallyConvex.Manifolds.CuspFinsler.eventually_outside_cusp, LocallyConvex.Manifolds.CuspFinsler.lengthStructure_not_finslerMetric}
  \leanok
  \uses{prop:mt-cusp-fibers, def:finsler-metric, def:finsler-length-structure}
  The field $\mu$ is a Finsler length structure but not a Finsler metric.
\end{theorem}

\begin{proof}
  \leanok
  It is not continuous on $\T\R^2$ at $(0, e_2)$: $\mu_{(t,0)}(e_2) = 2$ for $t \neq 0$, while
  $\mu_0(e_2) = 1$. Let now $\gamma$ be a $\Ck{1}$ curve. Away from the times where $\gamma$ passes
  through the origin, $\beta \circ \gamma$ is continuous. Suppose that $\gamma(t_0) = 0$. If
  $\dot\gamma_2(t_0) \neq 0$, then $|\gamma_2(t)| \ge c\,|t - t_0|$ while $\gamma_1(t)^2 =
  O(|t-t_0|^2)$, so near $t_0$ the curve stays outside the cusp and $\beta(\gamma(t)) = 0$. If
  $\dot\gamma_2(t_0) = 0$, then $0 \le \beta(\gamma)\,|\dot\gamma_2| \le |\dot\gamma_2| \to 0$.
\end{proof}

The example shows what a $\Ck{1}$ curve can and cannot see: through a point $p$ it only meets the
pairs $(p, v)$ that are compatible with the way it approaches $p$. The pairs $((1/n, 0), e_2)$,
which converge to $(0, e_2)$ in $\T\R^2$, are not of this kind.
